\begingroup
\shorthandoff{"<>"}
\setlength{\arrayrulewidth}{1pt}
\arrayrulecolor{vlepsbased}
\begin{longtable}{|p{15cm}|}
  \hline
  \rowcolor{maincolor} 
  \multicolumn{1}{|l|}{\textcolor{white}{\textbf{Caso de Uso}}} \\
  \hline
  \endfirsthead
  \multicolumn{1}{c}{{\tablename\ \thetable{} -- Continuación}} \\
  \hline
  \rowcolor{maincolor} 
  \multicolumn{1}{|l|}{\textcolor{white}{\textbf{Caso de Uso (Continuación)}}} \\
  \hline
  \endhead
  \hline \multicolumn{1}{|r|}{\textit{Continúa en la siguiente página...}} \\ \hline
  \endfoot
  \hline
  \caption{Caso de Uso CU11 — Eliminación permanente de sesión de chat.}
  \label{TB:CU11_ELIMINAR_SESION}
  \endlastfoot

  \rowcolor{ulepsbased}
  \begin{tabularx}{\linewidth}{@{}p{7.5cm}p{7.5cm}@{}}
    \textbf{Identificador:} CU11 & \textbf{Actor principal:} Usuario
  \end{tabularx} \\ \hline

  \textbf{Nombre:} Eliminación permanente de sesión de chat \\ \hline

  \rowcolor{ulepsbased}
  \textbf{Descripción:} \newline 
  El usuario decide eliminar permanentemente una sesión de chat completa de su historial, desencadenando el borrado en cascada de todos los mensajes asociados a la misma. \\ \hline

  \textbf{Precondiciones:} \newline
  1. El usuario ha iniciado sesión. \newline
  2. Existe al menos una sesión de chat. \\ \hline

  \rowcolor{ulepsbased}
  \textbf{Flujo principal:} \newline
  1. El usuario hace clic en el botón de eliminar (ícono de papelera) junto a una sesión en el panel lateral. \newline
  2. El frontend muestra un cuadro de diálogo de confirmación ("¿Estás seguro de que quieres eliminar esta sesión?"). \newline
  3. El usuario confirma la eliminación. \newline
  4. La interfaz envía una solicitud HTTP DELETE al backend. \newline
  5. El sistema verifica la propiedad y borra el registro de la sesión y, en cascada, todos sus mensajes asociados en la base de datos. \newline
  6. El sistema devuelve HTTP 204 (No Content). \newline
  7. El frontend remueve la sesión de la lista visual y, si el usuario estaba visualizando la sesión eliminada, limpia la pantalla y lo devuelve a la vista inicial (pantalla de creación de "Nueva conversación"), independientemente de si hay más sesiones en el historial. \\ \hline

  \textbf{Flujos alternativos:} \newline
  \textbf{FA-01:} Cancelación del borrado: \newline
  1. En el paso 2, el usuario pulsa "Cancelar". \newline
  2. Se cierra el cuadro de diálogo y no se envía ninguna solicitud al backend. \\ \hline

  \rowcolor{ulepsbased}
  \textbf{Postcondiciones:} \newline
  1. La sesión y todo su historial de mensajes desaparecen permanentemente del sistema. \\ \hline

\end{longtable}
\endgroup
